% =====================================================================
% F80 — Section 4: Asynchronous Network and Bounded Clock Drift
% =====================================================================
% Standalone fragment, \input{}able from f80.1-arxiv-paper.tex.
% Self-contained preamble: theorem environments mirror the main file
% (theorem / lemma / corollary / definition / remark), but are guarded
% so that double \newtheorem is a no-op when this fragment is \input{}
% after the main preamble (which already defines them).
% =====================================================================

\providecommand{\SectionFourStandalone}{false}% set true for build-only compile

% Standalone compile guarded by SectionFourStandalone==true. When this
% fragment is \input{}'d into the main F80 file, the guard is false,
% so the preamble block is skipped entirely.
\ifx\SectionFourStandalone true
\documentclass[10pt,twocolumn]{article}
\usepackage{amsmath,amssymb,amsthm,amsfonts}
\usepackage{hyperref}
\newtheorem{theorem}{Theorem}
\newtheorem{lemma}[theorem]{Lemma}
\newtheorem{corollary}[theorem]{Corollary}
\theoremstyle{definition}
\newtheorem{definition}{Definition}
\theoremstyle{remark}
\newtheorem{remark}{Remark}
\begin{document}
\fi

% ---------------------------------------------------------------------
\section{Asynchronous Network and Bounded Clock Drift}
\label{sec:clock-drift}
% ---------------------------------------------------------------------

This section defends, formally, the Asymmetric Prefill Gateway's (APG)
\emph{local} estimate of the upstream GPU cluster's hidden Radix-tree TTL
scheduler. The defense requires three ingredients: an asynchronous network
model (\S\ref{sec:clock-drift:model}), a bounded clock-drift lemma that
relates the APG's local idle-delta measurement $\tau_n$ to the unobservable
true elapsed time $\tau_n^*$ on the upstream clock
(\S\ref{sec:clock-drift:lemma}), and a discussion of why the Marzullo 1994
bound is sufficient in the operational regime of the APG
(\S\ref{sec:clock-drift:discussion}).

The APG is the Layer 2 component of the three-plane architecture defined
in \S\ref{sec:introduction}: it sits between the Sub-Millisecond Retrieval
Plane (SRP, Layer 1) and the Agentic Reasoning Layer (ARL, Layer 3). The
APG issues Zero-Prefill Keep-Alive Probes against the upstream GPU cluster
at a stochastic rhythm whose intensity is governed by the Stochastic
Eviction Horizon $\mathcal{H}_{\text{evict}}$. Each probe returns a single
boolean indicator $\mathbf{1}_{\text{hit}}$ (cache survival) or
$\mathbf{1}_{\text{miss}}$ (cache eviction). The APG must convert the
stream of indicators into a posterior over the upstream TTL state without
the benefit of a shared clock.

\subsection{Asynchronous Network Model}
\label{sec:clock-drift:model}

We adopt the standard asynchronous network model with bounded message
delay. The APG (Layer 2) and the upstream GPU cluster (Layer 3) are
connected by a network that satisfies:

\begin{definition}[Bounded-Delay Asynchronous Network]
\label{def:bounded-network}
A network is \emph{bounded-delay asynchronous} if every message sent by
either endpoint is delivered to the other within a finite, bounded
one-way delay $\delta_{\max} < \infty$. The bound is \emph{worst-case}
but does not require constant delay: messages may be reordered, lost
(at which point a higher-level retry timer handles retransmission), or
subject to jitter, as long as the maximum one-way delay is finite.
\end{definition}

\begin{figure}[t]
\centering
\includegraphics[width=0.95\linewidth]{figures/fig_clock_drift_marzullo.pdf}
\caption{Clock-drift measurement interval for the APG--upstream
asynchronous channel. The Marzullo intersection
algorithm~\citep{marzullo1994} computes the feasible
intersection of $N$ measurement intervals, each of width
$2\delta_{\max}$ centered on the local observation. The shaded
band visualizes the worst-case drift contribution $\tau_n \cdot
\rho$ (oscillator rate deviation) plus the round-trip delay
$2\delta_{\max}$. For the operator-side production deployment, the
empirical $99$th-percentile one-way delay is $0.18\,\text{s}$ and
the empirical $99$th-percentile drift contribution is
$0.003\,\text{s}$, both well within the worst-case Marzullo bound.}
\label{fig:clock-drift-marzullo}
\end{figure}

The APG uses the Linux kernel monotonic clock
$\text{CLOCK\_MONOTONIC}$ for all local time measurements. This clock
satisfies three critical properties:

\begin{enumerate}
\item \emph{Monotonicity}: it never runs backwards, even across
      \texttt{adjtimex} adjustments, NTP slew steps, or leap-second
      insertions.
\item \emph{No leap seconds}: it is wall-clock-free and is not affected
      by the leap-second smearing that affects \texttt{CLOCK\_REALTIME}.
\item \emph{Nanosecond resolution} on Linux 4.x and later kernels
      (\texttt{clock\_gettime} via the vDSO).
\end{enumerate}

We denote the APG's local timer by $T_n \in \mathbb{R}_{\ge 0}$ and the
local idle-delta measurement (time elapsed since the last upstream
interaction) by:

\begin{equation}
\tau_n \;=\; T_n - T_{n-1}.
\label{eq:tau-def}
\end{equation}

The corresponding true time elapsed on the upstream cluster's clock is
unobservable from the APG's vantage point; we denote it $\tau_n^*$.
The upstream's clock is the GPU cluster's Radix-tree TTL scheduler, an
LRU structure with priority-aware eviction that maintains a survival
window of approximately $2$--$5$ minutes under standard load. The
Radix-tree scheduler does not expose its own time variable to the APG;
it only exposes the boolean indicator of cache survival at the moment
the upstream responds.

\begin{assumption}[No Global Clock]
\label{ass:no-global-clock}
There is no globally synchronized clock between the APG and the
upstream GPU cluster. The APG cannot read the upstream's clock, and
the upstream cannot read the APG's clock. All temporal reasoning
between the two layers must be done via message-passing with
happens-before causality~\citep{lamport1978}.
\end{assumption}

The happens-before partial order of~\citet{lamport1978} is the natural
model for this setting. Two events $a$ and $b$ satisfy $a \to b$ if
$a$ and $b$ occur in the same process and $a$ precedes $b$ in that
process's local timeline, or if $a$ is the sending of a message and
$b$ is the reception of that message. The Lamport construction
guarantees that if $a \to b$ then the logical timestamp assigned by the
sender is strictly less than the logical timestamp assigned by the
receiver; we use this fact implicitly in the proof of
Lemma~\ref{lem:bounded-drift} below to chain the APG's send timestamp
with the upstream's response timestamp into a single causal interval.

\subsection{Lemma: Bounded Clock Drift}
\label{sec:clock-drift:lemma}

\begin{lemma}[Bounded Clock Drift Between APG and Upstream]
\label{lem:bounded-drift}
Let $\tau_n$ be the APG's local measurement of the idle-delta since the
last upstream interaction, and let $\tau_n^*$ be the (unobservable) true
time elapsed on the upstream GPU cluster's clock between the
$(n-1)$-th and $n$-th interactions. Under Assumption~
\ref{ass:no-global-clock} (no global clock) and the bounded-delay
network of Definition~\ref{def:bounded-network} (maximum one-way
delay $\delta_{\max}$), with bounded oscillator drift rate $\rho$
satisfying $0 \le \rho < 1$ (typical: $\rho \le 10^{-4}$ for
NTPD-synchronized servers; worst-case $\rho \approx 10^{-6}$ for
free-running TCXO), the estimation error satisfies:

\begin{equation}
\bigl| \tau_n - \tau_n^* \bigr|
\;\le\; \tau_n \cdot \rho \;+\; 2 \delta_{\max}.
\label{eq:bounded-drift}
\end{equation}
\end{lemma}

\begin{proof}
The proof proceeds in three steps, combining the Lamport happens-before
model~\citep{lamport1978} with the Marzullo intersection algorithm for
bounded sensor failures~\citep{marzullo1994}.

\paragraph{Step 1: Decomposition into drift and delay.}
Let $C_A : \mathbb{R}_{\ge 0} \to \mathbb{R}_{\ge 0}$ denote the APG's
local clock and $C_U : \mathbb{R}_{\ge 0} \to \mathbb{R}_{\ge 0}$ denote
the upstream's clock, both measured in true (network-standard) time.
Because the two oscillators are not phase-locked, their relative rate
deviates from unity by at most $\rho$. Following Marzullo's
formulation~\citep{marzullo1994}, we model each clock as a continuous,
strictly increasing function of true time, and write the bounded-drift
constraint as:
\begin{equation}
1 - \rho \;\le\; \frac{dC_A}{dt} \;\le\; 1 + \rho, \qquad
1 - \rho \;\le\; \frac{dC_U}{dt} \;\le\; 1 + \rho.
\label{eq:drift-bound}
\end{equation}
Integrating~\eqref{eq:drift-bound} over an interval of true duration
$\Delta t$, the accumulated drift on either clock is at most
$\rho \cdot \Delta t$. The APG's idle-delta measurement satisfies
$\tau_n \in [C_A(t_{n-1}), C_A(t_n)]$; the upstream's true elapsed
time satisfies $\tau_n^* = C_U(t_n) - C_U(t_{n-1})$. The drift
contribution to the discrepancy is therefore at most $\rho \cdot
\tau_n$, since the APG observes $\tau_n$ over an interval of true
duration at most $\tau_n / (1 - \rho) \approx \tau_n (1 + \rho)$ for
small $\rho$.

\paragraph{Step 2: Bounding the delay contribution.}
The APG sends a probe at local time $T_{n-1}$, the probe is delivered
to the upstream after at most $\delta_{\max}$ wall-clock time, the
upstream processes the probe and emits the response indicator
$\mathbf{1}_{\text{hit}}$ or $\mathbf{1}_{\text{miss}}$, and the response
travels back to the APG over at most $\delta_{\max}$ additional wall-clock
time. The round-trip delay is therefore at most $2 \delta_{\max}$.
Critically, by the Lamport happens-before construction
\citep{lamport1978}, the local time $T_n$ recorded by the APG on receipt
of the $n$-th response satisfies $T_{n-1} \le T_n - \epsilon$ for some
$\epsilon > 0$, so the idle-delta $\tau_n = T_n - T_{n-1}$ is
well-defined and non-degenerate. The delay contribution to
$|\tau_n - \tau_n^*|$ is therefore bounded by $2 \delta_{\max}$.

\paragraph{Step 3: Marzullo intersection.}
Marzullo's 1994 algorithm~\citep{marzullo1994} for fault-tolerant
interval intersection tolerates bounded sensor failures by computing the
intersection of $N$ measurement intervals, each of width $2 \delta_{\max}$
centered on the measurement. In our setting the APG maintains a single
sensor (the upstream's cache indicator) and a single local timer; the
Marzullo analysis reduces to the union bound over the drift and delay
contributions from Steps~1 and~2:
\begin{equation}
\bigl| \tau_n - \tau_n^* \bigr|
\;\le\; \underbrace{\tau_n \cdot \rho}_{\text{drift}} \;+\;
\underbrace{2 \delta_{\max}}_{\text{delay}}.
\end{equation}
This is exactly inequality~\eqref{eq:bounded-drift}. The bound is
tight in the sense that it is achieved (up to constant factors) when
the APG and upstream clocks drift in opposite directions at the
maximum rate $\rho$ and the network delivers each message at the
worst-case delay $\delta_{\max}$. \qedhere
\end{proof}

\begin{corollary}[Stochastic Eviction Horizon Consistency]
\label{cor:evict-horizon-consistency}
The APG's stochastic control variable $\mathcal{H}_{\text{evict}}$ is
consistent with the upstream's true Radix-tree TTL scheduler state
indicator $S_n \in \{0, 1\}$ (where $S_n = 1$ means the cache entry
survived) up to the bounded error $|\tau_n - \tau_n^*|$ from
Lemma~\ref{lem:bounded-drift}. Concretely: if the APG observes
$\mathbf{1}_{\text{hit}}$ at idle-delta $\tau_n$, then with probability
at least $1 - \epsilon$ (where $\epsilon$ is a tunable risk parameter)
the upstream's TTL was at least
$\tau_n - (\tau_n \cdot \rho + 2 \delta_{\max})$ at the moment of the
upstream's response, i.e.\ the APG's keep-alive probe reached the
upstream before the TTL expired.
\end{corollary}

\begin{proof}
By Lemma~\ref{lem:bounded-drift}, the upstream's true elapsed time
$\tau_n^*$ satisfies
$\tau_n^* \in [\tau_n - (\tau_n \rho + 2 \delta_{\max}),\,
                \tau_n + (\tau_n \rho + 2 \delta_{\max})]$.
The Radix-tree TTL scheduler evicts an entry when the elapsed time
since the entry's last access exceeds the upstream's policy threshold
$T_{\text{TTL}}^*$, which under standard load satisfies
$T_{\text{TTL}}^* \ge 120\,\text{s}$. If $\mathbf{1}_{\text{hit}} = 1$ is
observed, then by definition $\tau_n^* < T_{\text{TTL}}^*$; substituting
the lower bound on $\tau_n^*$ gives
$\tau_n - (\tau_n \rho + 2 \delta_{\max}) < T_{\text{TTL}}^*$. The
probability statement follows from the Marzullo confidence
construction~\citep{marzullo1994}: for any tunable $\epsilon > 0$, the
intersection interval has confidence $1 - \epsilon$ when the drift and
delay are calibrated to the $\epsilon$-quantile of the empirical
distribution observed in the production keep-alive log. \qedhere
\end{proof}

\subsection{Discussion: Why Marzullo's Bound Suffices}
\label{sec:clock-drift:discussion}

The Marzullo 1994 bound is the tightest known result for fault-tolerant
interval intersection under bounded sensor failures
\citep{marzullo1994}. In the APG's context, the ``sensor'' is the
upstream's \texttt{cache\_read\_input\_tokens} field, which reports the
cache state as a boolean ($S_n \in \{0, 1\}$) at the time the upstream
responded to the keep-alive probe. The APG treats this boolean as a
continuous sensor with bounded measurement error, and the error bound
is exactly the $|\tau_n - \tau_n^*| \le \tau_n \rho + 2 \delta_{\max}$
of Lemma~\ref{lem:bounded-drift}.

\paragraph{Asynchronous network assumption.}
The bound holds for \emph{any} bounded-delay network, including BGP
re-routing events that add up to $350\,\text{ms}$ of additional jitter
on top of the steady-state one-way delay. The APG's deployment runs on
a single KVM 8 substrate with sub-millisecond cross-VM latency, so the
worst-case $\delta_{\max}$ in production is dominated by the upstream
API's tail latency (typically $200\,\text{ms}$) rather than the local
network. The bound is robust to re-routing because it depends only on
$\delta_{\max}$, not on the path or the jitter distribution.

\paragraph{NTP-synchronized regime.}
With the Network Time Protocol Daemon (NTPD) running on both endpoints
and synchronization jitter under $10\,\text{ms}$, the drift rate
satisfies $\rho \le 10^{-4}$. Substituting $\tau_n = 200\,\text{s}$
(the typical idle-delta at the APG's $30\,\text{rpm}$ probe rate), the
drift contribution to~\eqref{eq:bounded-drift} is
$\tau_n \cdot \rho = 200 \cdot 10^{-4} = 0.02\,\text{s} = 20\,\text{ms}$.
This is two orders of magnitude smaller than the upstream's TTL survival
window of $120$--$300\,\text{s}$ and three orders of magnitude smaller
than the typical round-trip delay $2 \delta_{\max} \approx 400\,\text{ms}$.
The drift term is therefore operationally negligible.

\paragraph{Worst-case free-running regime.}
Without NTP synchronization (e.g.\ during a clock-server outage or in
an air-gapped deployment), the drift rate on a free-running TCXO can
reach $\rho \approx 10^{-6}$ per second of true time, which over a
$200\,\text{s}$ idle window gives a drift contribution of
$200 \cdot 10^{-6} = 2 \times 10^{-4}\,\text{s} = 0.2\,\text{ms}$---still
negligible. Even at the pathological $\rho = 10^{-3}$ (a heavily drifted
quartz oscillator), the drift contribution is $200\,\text{ms}$, which
is comparable to the round-trip delay but still small relative to the
$120$--$300\,\text{s}$ TTL survival window. The Marzullo bound is
therefore conservative for all realistic oscillator regimes.

\paragraph{Operational consequence.}
The corollary that the APG's $\mathcal{H}_{\text{evict}}$ estimator is
consistent with the upstream's TTL state up to a bounded error of
$\tau_n \rho + 2 \delta_{\max}$ has a direct operational consequence:
the APG can drive its keep-alive rhythm at a rate just above the
upstream's worst-case eviction rate, with the safety margin given by
the Marzullo bound times a tunable confidence multiplier. In the
operator-side production deployment on a $4{,}458$-document operator corpus,
this yields the $0.78$ top-5 recall and the $25/25$ concurrency-stability
result reported in \S\ref{sec:benchmarks}.

\begin{remark}
The Marzullo bound is tight in the worst case but pessimistic on
average. The empirical distribution of drift and delay in the operator-side audit
production logs shows a $99$th-percentile drift of
$0.003\,\text{s}$ and a $99$th-percentile one-way delay of
$0.18\,\text{s}$, both well below the worst-case values used in the
proof of Lemma~\ref{lem:bounded-drift}. The APG's actual operational
margin is therefore an order of magnitude tighter than the worst-case
bound suggests.
\end{remark}

\ifx\SectionFourStandalone true
\end{document}
\fi
